\documentclass[a4paper, 12pt]{article}
\usepackage[russian]{babel}
\usepackage[T2A]{fontenc}
\usepackage{indentfirst}
\usepackage[margin=1.5cm]{geometry}
\usepackage[hidelinks]{hyperref}
\usepackage{bookmark}
\usepackage{booktabs}
\usepackage{amsfonts}
\usepackage{amsmath}

\title{\textbf{Продолжение линейного функционала с сохранением нормы. Теорема Хана-Банаха.}}
\author{Берлет Максим гр. 3384 \\ Вариант 2}
\date{}


\begin{document}
\maketitle
\pagebreak

\tableofcontents
\pagebreak

\section{Условие задания.}
Пространство $X = \mathbb{R}^4$, $Y$ - подпростраство X; $Y = ker(k) = \{x \in X: (k, x) = 0\}; k \in X$.
$g \text{ - линейный функционал, } g \in Y^*: \forall y \in Y: g(y) = (g, y)$. 
Нужно построить продолжение функционала $g \in Y^*$ на все пространство $X$: $f \in X^*$ и $||f||_{X^*} = ||g||_{Y^*}$.

\section{Комментарии по условию.}
Фактически нам дано все пространство $X = \mathbb{R}^4$. Из него берется подпространство $Y$, которое фактически является гиперплоскостью с нормалью $\overrightarrow{k}$. На этой гиперплоскости задается функционал $g \in Y^*$ в стандартном базисе для $\mathbb{R}^4$.
По теореме Рисса-Фреше, так как наше пространство Гильбертово, то любой функционал можно записать в виде скалярного произведения: $g(y) = (g, y), y \in Y$.
И норма такого функционала $||g||_{Y^*} = A$, однако если попробовать вместо элементов из подпростраства(гиперплоскости), на котором определен функционал, взять элементы из изначального пространства, то обнаружится, что там норма функционала станет уже другой: $||g||_{X^*} = B \ne A$.
А это значит, что такой функционал на новом направлении(оно одно так как $X = ker(g) \oplus span\{x_0\}$ или $codim \ ker(g) = 1$), которого не было в подпространстве, растягивает больше.
Таким образом нужно найти такой функционал $f \in X^*$, что для всех элементов из $Y$ он будет совпадать с $g$, но во всем прострастве его норма будет равна норме исходного функционала.

\section{Теоретические выкладки.}

\section{Выполнение задания.}

\section{Выводы.}

\end{document}